% =====================================================================
% Clase -- Trigonometría 10° -- Trimestre I -- Semana 04
% Sesión 013 (mar 22 sep., 11:50, 50 min) · 014 (jue 24, 11:50, 50) ·
% 015 (jue 24, 13:40, 50) · 016 (vie 25, 10:10, 50)
% Tema: Razones trigonométricas
% Hilo conductor del trimestre I: «Medir lo inalcanzable»
% Tema 1 de la guía: Razones trigonométricas (tema:razones)
% Todas las medidas en métricas (decisión de la docente, 2026-09-19).
% Material del docente (no se entrega a los estudiantes).
% =====================================================================
\documentclass[11pt]{article}
\makeatletter\def\input@path{{../../../../../plantillas/estilo/}}\makeatother
\usepackage{mmcantillo}
\guiadelestudiante{../../guia-didactica/guia-periodo-I-medir-lo-inalcanzable}

\tipodocumento{Clase}
\asignatura{Trigonometría}
\titulo{Semana 4: las seis razones}
\grado{10°}
\periodo{I}

% DBA: matematicas grado 10 · DBA 4 — Comprende y utiliza funciones para modelar fenómenos
% periódicos y justifica las soluciones. Evidencias de esta semana: establecer las razones
% trigonométricas en el triángulo rectángulo y usarlas en los dos sentidos —de los lados al
% ángulo y del ángulo a los lados—, que es la base de toda la medición indirecta del hilo.

\begin{document}

\encabezado*

\begin{center}
{\renewcommand{\arraystretch}{1.3}
\begin{tabularx}{\linewidth}{@{}l l l X l@{}}
  \toprule
  \textbf{Sesión} & \textbf{Fecha} & \textbf{Min} & \textbf{Subtema} & \textbf{Entregables}\\
  \midrule
  013 & mar 22 sep., 11:50 & 50 & Razones dados dos lados (con Pitágoras) & ---\\
  014 & jue 24 sep., 11:50 & 50 & Dada una razón, hallar las otras & ---\\
  015 & jue 24 sep., 13:40 & 50 & Taller de cálculo de razones & Taller (5 puntos)\\
  016 & vie 25 sep., 10:10 & 50 & La calculadora científica & Tarea (5 puntos)\\
  \bottomrule
\end{tabularx}}
\end{center}

La semana en que la trigonometría deja de ser geometría y se vuelve cuenta. Tales midió la
pirámide con una sombra y una proporción; estas cuatro sesiones convierten esa proporción en
seis números con nombre propio, y enseñan a sacarlos en los dos sentidos.
Referencia: \refguia{tema:razones}.

% ---------------------------------------------------------------------
\seccion{Sesión 013 — martes 22 de septiembre · 11:50 · 50 min}

\textbf{Objetivo.} Calcular las seis razones trigonométricas de un ángulo agudo a partir de
dos lados del triángulo rectángulo, hallando el tercero con Pitágoras.

\begin{center}
{\renewcommand{\arraystretch}{1.3}
\begin{tabularx}{\linewidth}{@{}l l X@{}}
  \toprule
  \textbf{Momento} & \textbf{Min} & \textbf{Actividad}\\
  \midrule
  Enganche & 0--6 & Dos triángulos rectángulos semejantes en el tablero (catetos $3$--$4$ y
    $6$--$8$). Medir y comparar $\frac{\text{opuesto}}{\text{hipotenusa}}$ en los dos: da
    lo mismo. \textbf{La razón depende del ángulo, no del tamaño.} Ahí está la idea entera.\\
  Las seis & 6--22 & Nombrar seno, coseno y tangente respecto de un ángulo agudo
    \emph{concreto} (marcarlo siempre). Luego cosecante, secante y cotangente como sus
    inversas. Tabla en el tablero. Insistir en que «opuesto» y «adyacente» dependen de cuál
    de los dos ángulos agudos se mire.\\
  Con Pitágoras & 22--40 & Si faltan datos, primero el tercer lado. Ejemplo guiado: catetos
    $5$ y $12$, hipotenusa $13$, y las seis razones exactas como fracciones. Segundo
    ejemplo con el curso: $8$ y $15$ (hipotenusa $17$).\\
  Cierre & 40--50 & Que cada uno escriba las seis razones del \emph{otro} ángulo agudo del
    triángulo $5$--$12$--$13$ y compare: el seno de uno es el coseno del otro. Anotarlo, se
    usa en el trimestre II.\\
  \bottomrule
\end{tabularx}}
\end{center}

% ---------------------------------------------------------------------
\seccion{Sesión 014 — jueves 24 de septiembre · 11:50 · 50 min}

\textbf{Objetivo.} Dada una sola razón de un ángulo agudo, construir un triángulo auxiliar y
deducir las demás.

\begin{center}
{\renewcommand{\arraystretch}{1.3}
\begin{tabularx}{\linewidth}{@{}l l X@{}}
  \toprule
  \textbf{Momento} & \textbf{Min} & \textbf{Actividad}\\
  \midrule
  El problema & 0--8 & «Sé que $\tg \theta = \frac{3}{4}$ y nada más. ¿Cuánto vale
    $\sen \theta$?» Dejar que lo intenten un par de minutos antes de dar la técnica.\\
  Triángulo auxiliar & 8--26 & La técnica: la razón dice dos lados; se dibuja un triángulo
    con \emph{esos} números ($3$ opuesto, $4$ adyacente), se halla el tercero con Pitágoras
    ($5$) y se leen las demás razones. Comprobar con
    $\sen^2\theta + \cos^2\theta = 1$.\\
  ¿Y si cambio el tamaño? & 26--36 & Rehacerlo con $6$ y $8$: mismas razones. Conectar con
    el enganche del martes. Un ángulo, una razón — no importa el triángulo que se dibuje.\\
  Práctica & 36--50 & Dos o tres casos en el cuaderno, uno con $\sen$ dado y otro con
    $\cos$ dado, circulando por los puestos. Preparación directa del taller de la tarde.\\
  \bottomrule
\end{tabularx}}
\end{center}

% ---------------------------------------------------------------------
\seccion{Sesión 015 — jueves 24 de septiembre · 13:40 · 50 min}

\textbf{Objetivo.} Aplicar lo de las dos sesiones anteriores en el taller, con seguimiento
individual.

\begin{center}
{\renewcommand{\arraystretch}{1.3}
\begin{tabularx}{\linewidth}{@{}l l X@{}}
  \toprule
  \textbf{Momento} & \textbf{Min} & \textbf{Actividad}\\
  \midrule
  Repaso & 0--8 & Tres preguntas rápidas al aire: ¿qué es el seno? ¿de qué depende una
    razón? ¿qué se hace si solo dan una razón? Despejar dudas antes de empezar.\\
  Taller & 8--42 & Taller individual (5 puntos), sin calculadora: todas las respuestas son
    fracciones exactas. Circular. La pregunta 4 es un error ajeno y suele ser la que más
    discusión genera.\\
  Puesta en común & 42--50 & Recoger y resolver en el tablero solo la pregunta 4. Anuncio
    de la sesión 016: mañana sí entra la calculadora.\\
  \bottomrule
\end{tabularx}}
\end{center}

% ---------------------------------------------------------------------
\seccion{Sesión 016 — viernes 25 de septiembre · 10:10 · 50 min}

\textbf{Objetivo.} Usar la calculadora científica en modo grados para evaluar una razón y,
al revés, para hallar el ángulo a partir de la razón.

\textbf{Materiales.} \textbf{Calculadora científica} (avisarlo con un día de anticipación);
si faltan, trabajar en parejas.

\begin{center}
{\renewcommand{\arraystretch}{1.3}
\begin{tabularx}{\linewidth}{@{}l l X@{}}
  \toprule
  \textbf{Momento} & \textbf{Min} & \textbf{Actividad}\\
  \midrule
  \textbf{Modo grados} & 0--10 & \textbf{Lo primero, y lo que más errores causa en todo el
    año}: poner la calculadora en \textsc{deg} (no \textsc{rad} ni \textsc{grad}). Prueba
    de control en el tablero: $\sen 30^\circ$ debe dar exactamente $0{,}5$. Si no da,
    la calculadora está en otro modo. Que todos la hagan antes de seguir.\\
  Ida & 10--24 & Del ángulo a la razón: $\tg \num{54,4}^\circ$, $\cos \num{49,2}^\circ$.
    Convención de la clase: cuatro decimales para las razones. Y expresiones mixtas, como
    $\num{24,6}\cos\num{74,3}^\circ$, cuidando el orden de las teclas.\\
  Vuelta & 24--40 & De la razón al ángulo: las teclas $\sen^{-1}$, $\cos^{-1}$,
    $\tg^{-1}$ (casi siempre con \textsc{shift}). Ojo con el nombre: \emph{no} es el
    inverso multiplicativo. Ángulos a una décima de grado.\\
  Cierre & 40--50 & Dejar la tarea. Anuncio: con estas dos teclas ya se puede resolver un
    triángulo rectángulo completo — y con eso, medir la altura de un edificio desde la
    acera. Es el tema de la próxima semana.\\
  \bottomrule
\end{tabularx}}
\end{center}

\begin{nota}
\textbf{La prueba de los $30^\circ$.} Vale la pena volverla un ritual: cada vez que alguien
saque un resultado raro en lo que queda del año, lo primero es teclear $\sen 30^\circ$. Si no
da $0{,}5$, el problema no es el procedimiento sino el modo de la calculadora.
\end{nota}

\begin{nota}
\textbf{Unidades.} El módulo de las Guías de Apoyo trae los problemas de contexto en pies y
pulgadas. Según la decisión del 2026-09-19, el material de este curso va en métricas, así que
esos ejercicios quedaron fuera del taller y de la tarea; los que se usan son de cálculo puro o
ya vienen sin unidades.
\end{nota}

\end{document}
